% f02_canonical_tree variant d -- the canonical tree beside a legal-but-different
% association of the same eight operands, and why the roots are not the same word.
% Data: Definitions 1-3 of main.tex (sec:abi-tree); the non-associativity of
% IEEE-754 addition is stated at main.tex l.153; the canonical N=8 root word is
% the N=8 row of tables/invariance.tex.  The right-hand tree is the left-linear
% association -- a valid binary summation tree, not an SPRS object; no root word
% is claimed for it, only that it is a different one in general.
% Strategy: argue by COUNTEREXAMPLE.  a states the definition, b argues the
% addressing and c shows the construction is forced once the leaf placement is
% fixed.  All three answer "what is FBT(N)"; none answers "why does the ABI have
% to name one at all".  This variant does, by exhibiting a second tree that is
% perfectly legal arithmetic, computes the same real number, and returns a
% different 64-bit word -- so bit-exactness is a property of the chosen
% association, not of the summation.
% Colour: the canonical FBT is the compliant path (spBlue) closing on the pinned,
% confirmed root word (spGreen); the legal-but-different left-linear association
% is the violating alternative (spVerm) and its root is spVerm too. The two
% trees also differ in drawn shape and depth, so the contrast survives greyscale.
% Target: IEEE single column (3.5in = 8.9cm); drawn width approx 8.4cm.
\begin{tikzpicture}[
  x=1cm, y=1cm,
  inode/.style={draw=spBlue, circle, thick, inner sep=0pt, minimum size=4.2mm,
                font=\tiny, fill=spBlueF, text=spInk},
  altnode/.style={draw=spVerm, circle, thick, inner sep=0pt, minimum size=4.2mm,
                  font=\tiny, fill=spVermF, text=spInk},
  lnode/.style={draw=spSky!80!spInk, rounded corners=0.8pt, thick, inner sep=0pt,
                minimum size=4.0mm, font=\tiny, fill=spSkyF, text=spInk},
  rootn/.style={draw=spGreen, circle, line width=1.1pt, inner sep=0pt,
                minimum size=4.6mm, font=\tiny, fill=spGreenF, text=spInk},
  altroot/.style={draw=spVerm, circle, line width=1.1pt, inner sep=0pt,
                  minimum size=4.6mm, font=\tiny, fill=spVerm!22, text=spInk},
  edg/.style={draw=spBlue!70, thin},
  aedg/.style={draw=spVerm!75, thin, densely dashed},
  ttl/.style={font=\scriptsize\bfseries, anchor=north, align=center, text=spInk},
  sub/.style={font=\tiny, anchor=north, align=center, text=spInk},
  mid/.style={font=\tiny, anchor=north west, align=left, text width=2.30cm,
              inner sep=0pt, text=spInk},
]
\def\q{0.50}
\def\cx{6.45}

% =====================================================================
% LEFT -- FBT(8), the canonical association
% =====================================================================
\foreach \i in {0,...,7} {
  \pgfmathsetmacro{\xx}{0.05+\i*\q}
  \node[lnode] (C\i) at (\xx,0) {\i};
}
\foreach \n/\a/\b/\m in {3/0/1/0.5, 4/2/3/2.5, 5/4/5/4.5, 6/6/7/6.5} {
  \pgfmathsetmacro{\xx}{0.05+\m*\q}
  \node[inode] (CN\n) at (\xx,0.62) {\n};
  \draw[edg] (CN\n) -- (C\a); \draw[edg] (CN\n) -- (C\b);
}
\foreach \n/\a/\b/\m in {1/3/4/1.5, 2/5/6/5.5} {
  \pgfmathsetmacro{\xx}{0.05+\m*\q}
  \node[inode] (CN\n) at (\xx,1.24) {\n};
  \draw[edg] (CN\n) -- (CN\a); \draw[edg] (CN\n) -- (CN\b);
}
\node[rootn] (CN0) at ({0.05+3.5*\q},1.86) {0};
\draw[edg] (CN0) -- (CN1); \draw[edg] (CN0) -- (CN2);

\node[ttl, text=spBlue] at ({0.05+3.5*\q},5.05) {$\FBT(8)$\\ pinned by the ABI};
\node[sub] at ({0.05+3.5*\q},4.42)
  {depth $\lceil\log_2 8\rceil = 3$\\
   root $\Rcan(8,L)=$\\ \texttt{0x40E1948E978D4FDF}\\ (Table~I)};
\draw[spGreen, thin, -{Stealth[length=1.4mm]}]
  ({0.05+3.5*\q},3.30) -- ({0.05+3.5*\q},2.16);

% =====================================================================
% RIGHT -- the left-linear association over the same eight operands
% =====================================================================
\node[lnode] (A0) at (\cx,0) {0};
\node[lnode] (A1) at ({\cx+0.95},0) {1};
\node[altnode] (S1) at ({\cx+0.42},0.45) {};
\draw[aedg] (S1) -- (A0); \draw[aedg] (S1) -- (A1);
\foreach \k in {2,...,7} {
  \pgfmathsetmacro{\yy}{0.45*\k}
  \pgfmathsetmacro{\yl}{0.45*\k-0.23}
  \node[lnode] (A\k) at ({\cx+1.42},\yl) {\k};
  \ifnum\k=7 \node[altroot] (S\k) at ({\cx+0.42},\yy) {};
  \else      \node[altnode] (S\k) at ({\cx+0.42},\yy) {};\fi
  \pgfmathsetmacro{\kk}{\k-1}
  \draw[aedg] (S\k) -- (S\kk); \draw[aedg] (S\k) -- (A\k);
}
\node[ttl, text=spVerm] at ({\cx+0.86},5.05) {a legal\\ alternative};
\node[sub, text width=2.05cm] at ({\cx+0.86},4.42)
  {left-linear\\ depth $N{-}1=7$\\ root: a \emph{different}\\ FP64 word};
\draw[spVerm, thin, -{Stealth[length=1.4mm]}]
  ({\cx+0.42},3.86) -- ({\cx+0.42},3.40);
\node[font=\tiny, text=spVerm, anchor=south west, fill=white, inner sep=0.5pt] at ({\cx+0.62},3.24) {root};

% =====================================================================
% MIDDLE -- what the two trees share and where they part
% =====================================================================
\node[mid] at (3.95,4.72)
  {Same eight leaves, same order, same $+$ operator.\\[2pt]
   In $\mathbb{R}$ both roots are $\sum_k L[k]$.\\[2pt]
   In IEEE-754 they are not: addition is non-associative, so the
   \emph{shape} is part of the answer.};
\draw[spGrid, thin, dashed] (3.85,-0.10) -- (3.85,4.95);
\draw[spGrid, thin, dashed] (6.28,-0.10) -- (6.28,3.95);

% =====================================================================
% The consequence
% =====================================================================
\node[anchor=north west, font=\tiny, align=left, inner sep=0pt, text width=8.30cm,
      text=spInk] at (-0.10,-0.42)
  {\textbf{Why the choice matters.} Nothing about the hardware prefers the left
   tree: the right one is a correct summation, it is what a naive accumulator
   loop emits, and on a linear topology it is the cheaper communication pattern.
   The ABI pins the left one anyway, because a bit-exact root is only a
   \emph{comparison} if every run reduces in the same association --- otherwise
   two runs that disagree in the last ulp cannot be told apart from two runs one
   of which is broken. Pinning the tree is what converts the root word into a
   test. It costs nothing in freedom the ABI needs: topology, mapping,
   refinement and schedule all move nodes between PEs without moving a single
   node between parents.};
\end{tikzpicture}
